\section{Open Data Cube}

Основні переваги системи Data Cube:
\begin{itemize}
      \item Каталог великих геопросторових обсягів даних;
      \item API на основі Python для високопродуктивних запитів;
      \item Надає вченим та іншим користувачам легку можливість виконувати
	        дослідницький аналіз даних;
      \item Дозволяє масштабовану обробку даних континенту, що зберігається;
      \item Відстеження походження всіх даних, що містяться, щоб
	        забезпечити контроль якості та оновлення.
\end{itemize}

Мета ODC ~---~ збільшити роль супутникових даних %
наданням відкритого та доступного
дослідницького інструменту та сприяти %
розвитку сфери аналізу даних супутників.  Це
рішення має на меті підтримати ключові цілі, що %
включають в себе полегшення використання
даних та підтримувати глобальні пріоритетні %
міжнародні програми.  ODC повинен створити
"бренд", якому користувачі можуть довіряти, і %
він повинен сприяти позитивному
користувальницькому досвіду.  Це має стати %
можливим завдяки розробці спільноти ODC з
відкритим кодом, яка активно займається %
розробкою, ділиться алгоритмами та надає
підтримку один одному для вирішення проблем.


Як це рішення дозволяє задовольнити попит %
багатьох світових користувачів з обмеженими
ресурсами?  Масштабованість ~---~ одна з ключових %
проблем, з якою стикається ODC.  Хоча
початкові зусилля привели лише до декількох %
національних масштабів впровадження Data
Cube (наприклад, Австралія, Колумбія, Швейцарія), Є %
багато інших країн та глобальних
організацій, що мають високий інтерес до ODC.  %
Завдяки взаємодії з зацікавленими
організаціями та поточними користувачами, %
очікується, що кількість ODC збільшиться і
спільнота ODC процвітатиме з розширенням своїх %
можливостей.

Ядро ODC служить прошарком між супутниковими %
постачальниками даних та додатками.  Існує
набір інструментів з відкритим кодом, щоб %
допомогти вченим проводити дослідження,
використовуючи дані, якими керує ODC.  Популярні %
інструменти, що використовуються у
спільноті, яка використовує ODC Core як основу, %
включають:

\begin{itemize}
      \item \textbf{Командний рядок} інструмент, який %
використовують програмісти /
	    розробники для взаємодії з ODC;
      \item \textbf{Відкритий ODC} Візуальний та %
інтерактивний веб-додаток, який
      дозволяє
	    користувачам досліджувати свій інвентар %
наявних даних;
      \item \textbf{ODC статистика} Оптимізований засіб %
проведення розширеного аналізу в
	    системі ODC.  Цей інструмент орієнтований на %
вчених;
      \item \textbf{Веб-інтерфейс} ODC: веб-додаток, який %
дозволяє розробникам
      інтерактивно
	    демонструвати та візуалізувати вихід %
алгоритмів;
      \item \textbf{Jupyter Notebook} Такі файли дають %
інтерактивну можливість
	    візуалізувати аналізовані дані;
      \item \textbf{Веб-сервіси} відкритого %
геопросторового консорціуму: адаптери, які
	    можуть підключати додатки non-ODC до ODC
\end{itemize}


Усі вищеописані інструменти формують %
користувацьку структуру управління Open Data Cube.

У роботі був використаний ODC, який керується %
Інститутом космічних досліджень НАН
України та ДКА України.

\textbf{Причина використання ODC}

Причиною використання ODC у роботі стала %
розвиваюча спільнота, легкість та швидкість
завантаження знімків супутника родини Sentinel-2 і %
прикладні програми обробки з
використанням мови програмування Python, за %
допомогою яких відбувалось композитування,
візуалізація.  Але найбільша перевага ODC є %
масштабованість даних, що означає можливість
завантажувати і опрацьовувати як завгодно %
великі території.
